\begin{longtable}{ | p{5.2cm} | p{6.6cm} | p{2cm} | p{1.2cm} | }
\hline\rowcolor{lightgray}
\textbf{名称} & \textbf{描述} &  \textbf{地址} & \textbf{访问} \\ \hline
\multicolumn{4}{|l|}{\textbf{密钥寄存器}} \\ \hline

\hyperref[regdesc:AESKEYnREG]{AES\_KEY\_0\_REG}				&	AES 密钥资料寄存器 0				& 0x0000 & R/W \\ \hline
\hyperref[regdesc:AESKEYnREG]{AES\_KEY\_1\_REG}				&	AES 密钥资料寄存器 1				& 0x0004 & R/W \\ \hline
\hyperref[regdesc:AESKEYnREG]{AES\_KEY\_2\_REG}				&	AES 密钥资料寄存器 2				& 0x0008 & R/W \\ \hline
\hyperref[regdesc:AESKEYnREG]{AES\_KEY\_3\_REG}				&	AES 密钥资料寄存器 3				& 0x000C & R/W \\ \hline
\hyperref[regdesc:AESKEYnREG]{AES\_KEY\_4\_REG}				&	AES 密钥资料寄存器 4				& 0x0010 & R/W \\ \hline
\hyperref[regdesc:AESKEYnREG]{AES\_KEY\_5\_REG}				&	AES 密钥资料寄存器 5				& 0x0014 & R/W \\ \hline
\hyperref[regdesc:AESKEYnREG]{AES\_KEY\_6\_REG}				&	AES 密钥资料寄存器 6				& 0x0018 & R/W \\ \hline
\hyperref[regdesc:AESKEYnREG]{AES\_KEY\_7\_REG}				&	AES 密钥资料寄存器 7				& 0x001C & R/W \\ \hline

\multicolumn{4}{|l|}{\textbf{TEXT\_IN 寄存器}} \\ \hline

\hyperref[regdesc:AESTEXTINnREG]{AES\_TEXT\_IN\_0\_REG}			&	源数据资料寄存器 0					& 0x0020 & R/W \\ \hline
\hyperref[regdesc:AESTEXTINnREG]{AES\_TEXT\_IN\_1\_REG}			&	源数据资料寄存器 1					& 0x0024 & R/W \\ \hline
\hyperref[regdesc:AESTEXTINnREG]{AES\_TEXT\_IN\_2\_REG}			&	源数据资料寄存器 2					& 0x0028 & R/W \\ \hline
\hyperref[regdesc:AESTEXTINnREG]{AES\_TEXT\_IN\_3\_REG}			&	源数据资料寄存器 3					& 0x002C & R/W \\ \hline

\multicolumn{4}{|l|}{\textbf{TEXT\_OUT 寄存器}} \\ \hline

\hyperref[regdesc:AESTEXTOUTnREG]{AES\_TEXT\_OUT\_0\_REG}		&	结果数据资料寄存器 0					& 0x0030 & RO \\ \hline
\hyperref[regdesc:AESTEXTOUTnREG]{AES\_TEXT\_OUT\_1\_REG}		&	结果数据资料寄存器 1					& 0x0034 & RO \\ \hline
\hyperref[regdesc:AESTEXTOUTnREG]{AES\_TEXT\_OUT\_2\_REG}		&	结果数据资料寄存器 2					& 0x0038 & RO \\ \hline
\hyperref[regdesc:AESTEXTOUTnREG]{AES\_TEXT\_OUT\_3\_REG}		&	结果数据资料寄存器 3					& 0x003C & RO \\ \hline

\multicolumn{4}{|l|}{\textbf{控制或配置寄存器}} \\ \hline

\hyperref[regdesc:AESMODEREG]{AES\_MODE\_REG}				            & 配置密钥长度和加解密方向			    & 0x0040 & R/W \\ \hline
\hyperref[regdesc:AESDMAENABLEREG]{AES\_DMA\_ENABLE\_REG}		        &	配置 AES 加速器工作模式					                & 0x0090 & R/W \\ \hline
\hyperref[regdesc:AESBLOCKMODEREG]{AES\_BLOCK\_MODE\_REG}		        &	配置 DMA-AES 下的块运算模式			        & 0x0094 & R/W \\ \hline
\hyperref[regdesc:AESBLOCKNUMREG]{AES\_BLOCK\_NUM\_REG}		   	        &	配置块数量			        & 0x0098 & R/W \\ \hline
\hyperref[regdesc:AESINCSELREG]{AES\_INC\_SEL\_REG}			            &	配置标准增量函数               & 0x009C & R/W \\ \hline
\hyperref[regdesc:AESTRIGGERREG]{AES\_TRIGGER\_REG}			&	开始运算		& 0x0048 & WT \\ \hline
\hyperref[regdesc:AESCONTINUEOPREG]{AES\_CONTINUE\_OP\_REG}	&	继续运算	& 0x00A8 & WO \\ \hline
\hyperref[regdesc:AESDMAEXITREG]{AES\_DMA\_EXIT\_REG}		&	退出运算	& 0x00B8 & WO \\ \hline
\tagged{ESP32-P4-latest}{
\hyperref[regdesc:AESPSEUDOREG]{AES\_PSEUDO\_REG}				            & 配置伪轮次抗攻击功能			    & 0x00D0 & R/W \\ \hline
}

\multicolumn{4}{|l|}{\textbf{状态寄存器}} \\ \hline

\hyperref[regdesc:AESSTATEREG]{AES\_STATE\_REG}				&	运算状态			& 0x004C & RO \\ \hline

\multicolumn{4}{|l|}{\textbf{中断寄存器}} \\ \hline

\hyperref[regdesc:AESINTCLRREG]{AES\_INT\_CLR\_REG}			&	DMA-AES 中断清除	& 0x00AC & WT \\ \hline
\hyperref[regdesc:AESINTENAREG]{AES\_INT\_ENA\_REG}			&	DMA-AES 中断使能	& 0x00B0 & R/W \\ \hline

\end{longtable}
